\section*{Capítulo \ref{ch:b1:organic-redox} --- Oxidação e redução em química orgânica}

\begin{solution}{exo:b1:organic-redox:1}
Etanol: \ce{CH3} $-3$, \ce{CH2OH} $-1$. Etanal: \ce{CH3} $-3$, CHO $+1$.
Ácido etanoico: \ce{CH3} $-3$, COOH $+3$. Propanona: \ce{CH3} $-3$ (duas vezes),
C=O $+2$. Ácido metanoico: $+2$.
\end{solution}

\begin{solution}{exo:b1:organic-redox:2}
Propan-1-ol: propanal, depois ácido propanoico. Propan-2-ol: propanona.
2-Metilpropan-2-ol: nenhuma reação (a cor alaranjada permanece).
\end{solution}

\begin{solution}{exo:b1:organic-redox:3}
\ce{CH3COCH3 + 2H+ + 2e- -> CH3CH(OH)CH3};
\ce{CH3CHO + 2H+ + 2e- -> CH3CH2OH}.
\end{solution}

\begin{solution}{exo:b1:organic-redox:4}
\ce{NaBH4}: butan-1-ol, butan-2-ol, nenhuma reação, nenhuma reação (o
ácido é apenas desprotonado). \ce{LiAlH4}: butan-1-ol, butan-2-ol,
butan-1-ol e etanol, butan-1-ol.
\end{solution}

\begin{solution}{exo:b1:organic-redox:5}
\ce{5CH3CH(OH)CH3 + 2MnO4- + 6H+ -> 5CH3COCH3 + 2Mn^2+ + 8H2O}.
$n = 1.20/60.0 = \qty{0.0200}{mol}$; permanganato $\frac25 \times 0.0200 =
\qty{8.0e-3}{mol}$, \qty{400}{mL} da solução a \qty{0.020}{mol/L}.
\end{solution}

\begin{solution}{exo:b1:organic-redox:6}
Aqueça o álcool com o \omterm{def:b1:nernst:couple}{oxidante} em um balão montado para destilação, a uma
temperatura entre as duas temperaturas de ebulição: o propanal
(\qty{48}{\celsius}) destila à medida que se forma e escapa à oxidação
posterior; o propan-1-ol (\qty{97}{\celsius}) fica no balão.
\end{solution}

\begin{solution}{exo:b1:organic-redox:7}
Um par B--H de \ce{BH4-} ataca o carbono carbonílico, com o par $\pi$ da
C=O indo para o oxigênio; o \omterm{def:b1:alcohol-activation:alkoxide}{alcóxido} capta um próton do etanol. O
hidrogênio no carbono do produto (\ce{CH3CH2OH}, um dos dois em C1) vem
do reagente; o do oxigênio, do solvente.
\end{solution}

\begin{solution}{exo:b1:organic-redox:8}
O butan-2-ol é \omterm{def:b1:stereochemistry-in-depth:stereocentre}{quiral}, mas a cetona plana é atacada igualmente pelas duas
faces: uma \omterm{def:b1:stereochemistry-in-depth:enantiomers}{mistura racêmica}, opticamente inativa.
\end{solution}

\begin{solution}{exo:b1:organic-redox:9}
\ce{3CH3CH2OH + 2Cr2O7^2- + 16H+ -> 3CH3COOH + 4Cr^3+ + 11H2O}: o dicromato
alaranjado (cromo $+$VI) se torna cromo(III) verde.
\end{solution}

\begin{solution}{exo:b1:organic-redox:10}
2-Metilbutan-2-ol: etanal + \ce{CH3CH2MgBr} dá butan-2-ol (uma adição,
sem oxirredução); oxidação a butanona; \ce{CH3MgBr} sobre a butanona, e
depois o tratamento. Uma oxidação. Pentan-3-ol: propanal + \ce{CH3CH2MgBr}
em uma etapa, sem oxirredução.
\end{solution}

\begin{solution}{exo:b1:organic-redox:11}
1. Proteja a cetona como \omterm{def:b1:acetals-protection:hemiacetal}{acetal} cíclico (etano-1,2-diol, ácido,
Dean--Stark). 2. \ce{LiAlH4} em éter seco reduz o éster ao álcool primário
(e metanol). 3. O ácido aquoso remove o \omterm{def:b1:acetals-protection:hemiacetal}{acetal}:
5-hidroxipentan-2-ona, \ce{CH3CO(CH2)3OH}.
\end{solution}

\begin{solution}{exo:b1:organic-redox:12}
\ce{RCH2OH}: carbono carbinólico em $-1$; \ce{RCOOH}: $+3$; quatro elétrons, como em
\ce{RCH2OH + H2O -> RCOOH + 4H+ + 4e-}. Do metano ($-4$) ao dióxido de
carbono ($+4$): oito elétrons, \ce{CH4 + 2H2O -> CO2 + 8H+ + 8e-}.
\end{solution}

\begin{solution}{pb:b1:organic-redox:1}
\textbf{1.} C1 (que leva o OH): 0; os cinco \ce{CH2}: $-2$.
\textbf{2.} C1 (C=O): $+2$; os cinco \ce{CH2}: $-2$.
\textbf{3.} Os dois carbonos COOH: $+3$; os quatro \ce{CH2}: $-2$.
\textbf{4.} C1 ($0 \to +3$) e seu vizinho C6 ($-2 \to +3$), que se torna
o segundo carbono carboxílico; a ligação C1--C6 é rompida.
\textbf{5.} $3 + 5 = 8$ elétrons.
\textbf{6.} O ácido nítrico é forte o bastante para romper a ligação C--C
vizinha à carbonila, o que \omterm{def:b1:nernst:couple}{oxidantes} mais brandos não fazem.
\textbf{7.} \ce{C6H12O + 3H2O -> C6H10O4 + 8H+ + 8e-}.
\textbf{8.} $+$V e $+$I.
\textbf{9.} \ce{2HNO3 + 8H+ + 8e- -> N2O + 5H2O}.
\textbf{10.} Somando, oito elétrons, \ce{8H+} e três moléculas de água se
cancelam: \ce{C6H12O + 2HNO3 -> C6H10O4 + N2O + 2H2O}.
\textbf{11.} 8.
\textbf{12.} $10^6/146.0 = \qty{6.85e3}{mol}$ de ácido; $2 \times 6.85 \times
10^3 \times 63.0 = \qty{8.6e5}{g}$, \qty{0.86}{t} de ácido nítrico.
\textbf{13.} \ce{H2O2 + 2H+ + 2e- -> 2H2O}.
\textbf{14.} \ce{C6H12O + 4H2O2 -> C6H10O4 + 5H2O}.
\textbf{15.} $4 \times 6.85 \times 10^3 \times 34.0 = \qty{9.3e5}{g}$,
\qty{0.93}{t}.
\textbf{16.} A água: nenhum óxido de nitrogênio.
\textbf{17.} Uma constante de equilíbrio grande nada diz sobre a
velocidade; \ce{H2O2} reage lentamente sem \omterm{def:b1:elementary-steps:catalyst}{catalisador}, como na sua
própria decomposição.
\textbf{18.} Ele reduz a ciclo-hexanona de volta a ciclo-hexanol:
\ce{4C6H10O + BH4- -> B(OC6H11)4-}, depois
\ce{B(OC6H11)4- + 4H2O -> 4C6H11OH + B(OH)4-}.
\textbf{19.} \qty{6.85e3}{mol}.
\textbf{20.} \qty{6.85e3}{mol} de \ce{N2O} (um por ácido).
\textbf{21.} $6.85 \times 10^3/0.95 \times 100.0 = \qty{7.2e5}{g}$, \qty{0.72}{t}.
\textbf{22.} $6.85 \times 10^3 \times 44.0 = \qty{3.0e5}{g}$: \textbf{\qty{0.30}{t}
de \ce{N2O} por tonelada de ácido adípico}.
\end{solution}
